\section{Conclusion}

The Guestbook application is running on OpenShift with all four services deployed and
communicating correctly. The frontend is accessible over HTTPS through an automatically
assigned Route hostname, the backend connects to both PostgreSQL and Redis, and
PostgreSQL uses persistent storage so entries survive pod restarts. The container images
are built automatically by GitHub Actions on every push to main, tagged with the commit
SHA, scanned for vulnerabilities, and published to GHCR. The infrastructure manifests
live in a separate repository as the assignment required.

The ArgoCD integration was not completed due to the Sandbox restrictions described in
section 5. Everything that was needed on the application and infrastructure side is in
place. The gap is purely a platform permission issue that would not exist on the shared
cluster or on any cluster where the user has cluster-admin access.

Looking at the project as a whole, I think the main thing it demonstrated is that
cloud-native deployment is not just about writing the right YAML. A lot of the work is
understanding the security model of the platform you are deploying to: what users and
containers are allowed to do, how networking and DNS work, and where secrets should live.
These are things that are much clearer to me now than they were at the start of the
assignment.
